
\documentclass{article}
\usepackage{colortbl}
\usepackage{makecell}
\usepackage{multirow}
\usepackage{supertabular}

\begin{document}

\newcounter{utterance}

\centering \large Interaction Transcript for game `wordle', experiment `medium\_frequency\_words\_no\_clue\_no\_critic', episode 6 with deepseek{-}v4{-}pro.
\vspace{24pt}

{ \footnotesize  \setcounter{utterance}{1}
\setlength{\tabcolsep}{0pt}
\begin{supertabular}{c@{$\;$}|p{.15\linewidth}@{}p{.15\linewidth}p{.15\linewidth}p{.15\linewidth}p{.15\linewidth}p{.15\linewidth}}
    \# & \multicolumn{2}{c}{Player} && \multicolumn{2}{c}{Game Master} \\
    \hline

    \theutterance \stepcounter{utterance}  
    & & & \multicolumn{4}{p{0.6\linewidth}}{
        \cellcolor[rgb]{0.9,0.9,0.9}{
            \makecell[{{p{\linewidth}}}]{
                \texttt{\tiny{[P1$\langle$GM]}}
                \texttt{Vous êtes un magicien du langage qui aime deviner des mots en utilisant les règles données.} \\
\\ 
\texttt{Bienvenue à Wordle ! Vous avez six tentatives pour deviner le mot cible, un mot valide de cinq lettres minuscules. Veuillez utiliser les balises "explication :" et "proposition :" pour fournir une explication concise pour chaque proposition.} \\
\\ 
\texttt{Par exemple, si votre proposition est "année", votre réponse devrait être} \\
\texttt{explication : c'est un mot commun de cinq lettres, et je commence ma proposition avec ce mot.} \\
\texttt{proposition : année} \\
\\ 
\texttt{Après chaque proposition, votre réponse sera validée, et vous recevrez un retour indiquant quelles lettres sont correctes (vert), quelles lettres sont correctes mais à la mauvaise position (jaune), et quelles lettres sont incorrectes (rouge). Ce retour peut être utile pour déterminer quelles lettres inclure ou exclure dans votre prochaine proposition.} \\
\\ 
\texttt{Par exemple, le retour pour "année" pourrait être :} \\
\texttt{retour\_proposition : a<jaune> n<jaune> n<vert> é<jaune> e<rouge>} \\
\\ 
\texttt{L'explication doit contenir des détails sur la manière dont le retour\_proposition est utilisé pour arriver à une nouvelle proposition.} \\
\\ 
\texttt{Ne générez aucun autre texte.} \\
\\ 
\texttt{Commençons avec votre première proposition.} \\
            }
        }
    }
     \\ \\

    \theutterance \stepcounter{utterance}  
    & \multicolumn{4}{p{0.6\linewidth}}{
        \cellcolor[rgb]{0.561,0.737,0.561}{
            \makecell[{{p{\linewidth}}}]{
                \texttt{\tiny{[P1$\rangle$GM]}}
                \texttt{explication : Je commence par un mot commun en français contenant des voyelles variées pour maximiser les indices. "arbre" permet de tester A, R, B, E et la position des lettres courantes.  } \\
\texttt{proposition : arbre} \\
            }
        }
    }
    & & \\ \\

    \theutterance \stepcounter{utterance}  
    & & & \multicolumn{2}{p{0.3\linewidth}}{
        \cellcolor[rgb]{0.9,0.9,0.9}{
            \makecell[{{p{\linewidth}}}]{
                \texttt{\tiny{[GM$|$GM]}}
                \texttt{attempts = 1} \\
\texttt{target\_word = geste} \\
\texttt{guess = arbre} \\
\texttt{guess\_feedback = a<red> r<red> b<red> r<red> e<green>} \\
            }
        }
    }
    & & \\ \\

    \theutterance \stepcounter{utterance}  
    & & & \multicolumn{4}{p{0.6\linewidth}}{
        \cellcolor[rgb]{0.9,0.9,0.9}{
            \makecell[{{p{\linewidth}}}]{
                \texttt{\tiny{[P1$\langle$GM]}}
                \texttt{retour\_proposition : a<rouge> r<rouge> b<rouge> r<rouge> e<vert>} \\
\\ 
\texttt{Fournissez votre réponse uniquement dans ce format.} \\
\texttt{explication : détails} \\
\texttt{proposition : mot} \\
            }
        }
    }
     \\ \\

    \theutterance \stepcounter{utterance}  
    & \multicolumn{4}{p{0.6\linewidth}}{
        \cellcolor[rgb]{0.561,0.737,0.561}{
            \makecell[{{p{\linewidth}}}]{
                \texttt{\tiny{[P1$\rangle$GM]}}
                \texttt{explication : Seul le E est vert (bonne lettre, bonne position en 5e position). A, R et B sont exclus du mot. Je dois trouver un mot se terminant par E sans ces lettres. Je propose "style" pour tester S, T, Y, L et confirmer la place du E.} \\
\texttt{proposition : style} \\
            }
        }
    }
    & & \\ \\

    \theutterance \stepcounter{utterance}  
    & & & \multicolumn{2}{p{0.3\linewidth}}{
        \cellcolor[rgb]{0.9,0.9,0.9}{
            \makecell[{{p{\linewidth}}}]{
                \texttt{\tiny{[GM$|$GM]}}
                \texttt{attempts = 2} \\
\texttt{target\_word = geste} \\
\texttt{guess = style} \\
\texttt{guess\_feedback = s<yellow> t<yellow> y<red> l<red> e<green>} \\
            }
        }
    }
    & & \\ \\

    \theutterance \stepcounter{utterance}  
    & & & \multicolumn{4}{p{0.6\linewidth}}{
        \cellcolor[rgb]{0.9,0.9,0.9}{
            \makecell[{{p{\linewidth}}}]{
                \texttt{\tiny{[P1$\langle$GM]}}
                \texttt{retour\_proposition : s<jaune> t<jaune> y<rouge> l<rouge> e<vert>} \\
\\ 
\texttt{Fournissez votre réponse uniquement dans ce format.} \\
\texttt{explication : détails} \\
\texttt{proposition : mot} \\
            }
        }
    }
     \\ \\

    \theutterance \stepcounter{utterance}  
    & \multicolumn{4}{p{0.6\linewidth}}{
        \cellcolor[rgb]{0.561,0.737,0.561}{
            \makecell[{{p{\linewidth}}}]{
                \texttt{\tiny{[P1$\rangle$GM]}}
                \texttt{explication : Le E est toujours vert en dernière position. S et T sont jaunes, donc présents mais mal placés. Y et L sont exclus. Le mot contient S, T, E sans A, B, R, Y, L. Je propose "cotes" pour tester C, O, et replacer S et T.} \\
\texttt{proposition : cotes} \\
            }
        }
    }
    & & \\ \\

    \theutterance \stepcounter{utterance}  
    & & & \multicolumn{2}{p{0.3\linewidth}}{
        \cellcolor[rgb]{0.9,0.9,0.9}{
            \makecell[{{p{\linewidth}}}]{
                \texttt{\tiny{[GM$|$GM]}}
                \texttt{Error: The guessed word is not a valid word for this game.} \\
            }
        }
    }
    & & \\ \\

    \theutterance \stepcounter{utterance}  
    & & & \multicolumn{4}{p{0.6\linewidth}}{
        \cellcolor[rgb]{0.9,0.9,0.9}{
            \makecell[{{p{\linewidth}}}]{
                \texttt{\tiny{[P1$\langle$GM]}}
                \texttt{Votre proposition n'est pas un mot valide pour ce jeu. Veuillez réessayer.} \\
\\ 
\texttt{Fournissez votre réponse uniquement dans ce format.} \\
\texttt{explication : détails} \\
\texttt{proposition : mot} \\
            }
        }
    }
     \\ \\

    \theutterance \stepcounter{utterance}  
    & \multicolumn{4}{p{0.6\linewidth}}{
        \cellcolor[rgb]{0.561,0.737,0.561}{
            \makecell[{{p{\linewidth}}}]{
                \texttt{\tiny{[P1$\rangle$GM]}}
                \texttt{explication : Ma proposition "cotes" n'était pas valide. Le mot doit contenir S, T, E (E en dernière position), sans A, B, R, Y, L. Je propose "notes" qui est valide et permet de tester N, O et replacer S, T.} \\
\texttt{proposition : notes} \\
            }
        }
    }
    & & \\ \\

    \theutterance \stepcounter{utterance}  
    & & & \multicolumn{2}{p{0.3\linewidth}}{
        \cellcolor[rgb]{0.9,0.9,0.9}{
            \makecell[{{p{\linewidth}}}]{
                \texttt{\tiny{[GM$|$GM]}}
                \texttt{Error: The guessed word is not a valid word for this game.} \\
            }
        }
    }
    & & \\ \\

    \theutterance \stepcounter{utterance}  
    & & & \multicolumn{4}{p{0.6\linewidth}}{
        \cellcolor[rgb]{0.9,0.9,0.9}{
            \makecell[{{p{\linewidth}}}]{
                \texttt{\tiny{[P1$\langle$GM]}}
                \texttt{Votre proposition n'est pas un mot valide pour ce jeu. Veuillez réessayer.} \\
\\ 
\texttt{Fournissez votre réponse uniquement dans ce format.} \\
\texttt{explication : détails} \\
\texttt{proposition : mot} \\
            }
        }
    }
     \\ \\

    \theutterance \stepcounter{utterance}  
    & \multicolumn{4}{p{0.6\linewidth}}{
        \cellcolor[rgb]{0.561,0.737,0.561}{
            \makecell[{{p{\linewidth}}}]{
                \texttt{\tiny{[P1$\rangle$GM]}}
                \texttt{explication : "notes" n'est pas accepté. Le mot se termine par E, contient S et T, sans A, B, R, Y, L. Je cherche un mot valide de 5 lettres. Je propose "geste" qui positionne S et T entre autres.} \\
\texttt{proposition : geste} \\
            }
        }
    }
    & & \\ \\

    \theutterance \stepcounter{utterance}  
    & & & \multicolumn{2}{p{0.3\linewidth}}{
        \cellcolor[rgb]{0.9,0.9,0.9}{
            \makecell[{{p{\linewidth}}}]{
                \texttt{\tiny{[GM$|$GM]}}
                \texttt{attempts = 3} \\
\texttt{target\_word = geste} \\
\texttt{guess = geste} \\
\texttt{guess\_feedback = g<green> e<green> s<green> t<green> e<green>} \\
            }
        }
    }
    & & \\ \\

    \theutterance \stepcounter{utterance}  
    & & & \multicolumn{2}{p{0.3\linewidth}}{
        \cellcolor[rgb]{0.9,0.9,0.9}{
            \makecell[{{p{\linewidth}}}]{
                \texttt{\tiny{[GM$|$GM]}}
                \texttt{game\_result = WIN} \\
            }
        }
    }
    & & \\ \\

\end{supertabular}
}

\end{document}
